\begin{center}
  \Large
  \textbf{KATA PENGANTAR}
\end{center}

\addcontentsline{toc}{chapter}{KATA PENGANTAR}

\vspace{2ex}

Puji syukur penulis panjatkan ke hadirat Allah SWT atas segala rahmat,
karunia, dan petunjuk-Nya, sehingga penulis dapat menyelesaikan
penyusunan tugas akhir yang berjudul ``\tatitle{}'' dengan baik dan
tepat waktu. Tugas akhir ini disusun sebagai salah satu syarat untuk
memperoleh gelar Sarjana Teknik pada Departemen Teknik Elektro,
Fakultas Teknologi Elektro dan Informatika Cerdas, Institut Teknologi
Sepuluh Nopember.

Dalam proses pengerjaan dan penyusunan tugas akhir ini, penulis
menyadari bahwa kelancaran dan keberhasilannya tidak lepas dari
bantuan, bimbingan, dukungan, dan semangat dari berbagai pihak.
Oleh karena itu, dengan segala kerendahan hati, penulis ingin
menyampaikan terima kasih yang sebesar-besarnya kepada:

\begin{enumerate}
\item Allah SWT, atas segala rahmat, karunia, dan kemudahan yang
      senantiasa menyertai langkah penulis dalam menyelesaikan
      tugas akhir ini.

\item Kedua orang tua tercinta, Ayah dan Bunda, Adik dan Nenek atas doa yang tak pernah putus, dukungan moral
      maupun material, serta kasih sayang yang tak ternilai.

\item Bapak Dr. Rudy Dikairono, S.T., M.T., selaku dosen pembimbing
      pertama sekaligus pembimbing Tim IRIS, dan Bapak Muhammad Attamimi, B.Eng., M.Eng., Ph.D.,
      selaku dosen pembimbing kedua, yang telah meluangkan waktu,
      memberikan bimbingan, arahan, serta motivasi selama proses
      pengerjaan tugas akhir.

\item \begin{sloppypar}Azzam Wildan Maulana, Muhammad Hafidzh, Figo Arzaki Maulana, Faizzul Haq Alma'anizam Zami, Hafiedz Rosyada, dan Muhammad Raihan Ramadhan,
      selaku tim riset mobil Omoda ITS, yang bekerja sama dengan penulis dalam pengumpulan data, pengujian, dan
      pengembangan sistem navigasi mobil otonom.\end{sloppypar}

\item Teman-teman Tim IRIS sebagai teman-teman terdekat penulis yang
      senantiasa menjadi tempat berkembang dan belajar selama awal perkuliahan hingga
      akhir perkuliahan.

\item Teman-teman Lab AJ403, Lab Elektronika yang telah memberikan dukungan, motivasi, dan semangat selama proses pengerjaan tugas akhir.

\item Seluruh pihak lain yang tidak dapat disebutkan satu per satu,
      namun telah memberikan bantuan dan dukungan hingga laporan ini
      dapat diselesaikan.
\end{enumerate}

Penulis menyadari bahwa tugas akhir ini masih memiliki kekurangan.
Oleh karena itu, penulis membuka diri untuk menerima saran dan kritik
yang membangun demi perbaikan di masa yang akan datang. Besar harapan
penulis, semoga tugas akhir ini dapat memberikan manfaat bagi semua
pihak yang membutuhkan dan bagi pengembangan ilmu pengetahuan.

\vspace{2ex}

\begin{flushright}
Surabaya, \MONTH{} \the\year

\vspace{2ex}

\name{}
\end{flushright}
